\section{\centering Solutions}

\subsection{Solution to T/F Questions}
		\begin{table}[!h]
			\centering
			\begin{tabular}{c@{\hspace{2.5em}}cp{0.8\textwidth}}
				true & false & \\
				\specialrule{\lightrulewidth}{0.25em}{1.25em}
			         \bubble & \dbubble & For any square matrix \(AB\), \(\det(AB) = \det A \cdot \det B\).\\[0.75em]
				 \dbubble & \bubble & If \(B\) is a \(n \times n\) matrix and \(\colspace B \ne \mathbb{R}^n\), then for any \(n \times n\) matrix \(A\), we
								have that \(\det(AB) = \det(B)\).\\[1.75em]
				\bubble & \dbubble & Given that \(A\) is a square matrix and we obtain a matrix \(B\) by scaling a row of \(A\) by 2, then 
							      \(\det A = 2 \cdot \det B\). \\[1.75em]
				\dbubble & \bubble & If \(\det(A^{2023}) = 5\), then \(A\) is non-singular. \\[0.75em]
				\bubble & \dbubble & While the determinant of a lower triangular matrix is the product of it's diagonals, this is not the case for upper
							      triangular matrices. \\[1.75em]
				\dbubble & \bubble & For a square matrix \(A\) that has an \(LU\)-factorization \(A = LU\), \(\det(A) = \det(U)\).\\[0.75em]
				\dbubble & \bubble & Scaling a column of a square matrix by a factor of 2 doubles the determinant of the matrix.\\[0.75em]
				\dbubble & \bubble & If \(A\) is an invertible matrix, then \(\displaystyle \det(A) = \frac{1}{\det(A^{-1})}\).\\[0.75em]
			\bottomrule
			\end{tabular}
		\end{table}

\subsection{Solution to Short Answer Questions}
\begin{enumerate}
	\item 36.
	\item 0.
	\item 2.
	\item \(-24/2^n\).
	\item \(k = 18\). You do this problem through setting the determinant to 0 and finding the value \(k\) that does this.
	\item First, a picture.
		\begin{figure}[!h]
			\centering
			\begin{tikzpicture}
				\begin{axis}[axis lines=left, xmin=0, xmax=5, ymin=0, ymax=5,grid=major]
					\draw[fill=teal, opacity=0.25] (axis cs:0,0) -- (axis cs:1,2) -- (axis cs:4,3) -- (axis cs:3,1)  -- (axis cs:0,0);
					\addplot[->, semithick, domain=0:1, color=teal] {2*x} node[above, fill=white] {\(\begin{bmatrix}1 \\ 2\end{bmatrix}\)};
					\addplot[->, semithick, domain=0:3, color=teal] {1/3*x} node[right, fill=white] {\(\begin{bmatrix}3 \\ 1\end{bmatrix}\)};
					\addplot[->, dashed, domain=0:4] {3/4 * x} node[right, fill=white] {\(\begin{bmatrix}4 \\ 3\end{bmatrix}\)};
				\end{axis}
			\end{tikzpicture}
			\caption{Pictorial representation of the problem.}
		\end{figure}
		\begin{enumerate}
			\item The signed area of this shape is the determinant of the matrix \(\begin{bmatrix}1 & 3 \\ 2 & 1\end{bmatrix}\). We find that the
			 	 determinant is \(-5\). The area is therefore 5.
			\item We \emph{could} apply the transformation on each point and find the determinant of a new matrix, but a simpler way is to rely on the formula
				\[V_T = V_S \cdot \lvert \det A \rvert.\]
				Applying the above formula, we get that \(V_T\) (the volume of the transformed shape) is equal to \(V_S\) (the original volume) times
				\(\lvert \det A \rvert = 3.\) So, \(V_T = 15\).
		\end{enumerate}
\end{enumerate}
